\chapter{相关实验现象与模型检验}
\section{MRX四极Hall场：直接测到什么，推断了什么}
\citet{Ren2005}在MRX中测得重联平面外的四极磁场，并比较氢和氘放电、不同碰撞性条件。局地高分辨磁探针阵列间距1.25 mm，误差小于约10 G；二维图使用更粗的空间扫描与重复放电平均。重联场约100 G时，高分辨测得Hall场可达50--60 G，粗二维扫描则约20--40 G，说明仪器分辨率本身影响峰值。

实验四极场宽度与离子惯性长度同阶、长度约数个 $d_i$；随着电流片尺度与平均自由程的比值增大，Hall场相对重联场减弱。它是离子与电子解耦电流系统的直接证据，而不是仅从速率反推Hall效应。

但原文给出的约20 km/s电子流是根据磁场与Ampère关系推断；Hall电场约140 V/m与测得重联场约100 V/m的比较在上游区域完成，且某个磁分量有明显不确定性。该文没有完成电子X点全部欧姆定律项的局地闭合。四极场与高频扰动可共存，因此也不能由它排除其他非理想过程。

建议对照原文图2与图4：前者看结构及测量布局，后者看碰撞性变化。学习时应把“直接测量的磁场”“由它推断的电流／速度”和“对断线机制的解释”写成三列。

\section{MRX导引场：重联减慢同时包含局地与全局效应}
\citet{GuideExp2012}扫描导引场大小与极性，发现重联速率降低，同时面内Hall相关结构改变。装置参数中片半宽约2 cm，$d_i$ 约5 cm，声回旋半径约2.5 cm，几个尺度彼此接近，不能套用无限分离的渐近模型。

该文的Hall相关测量选在离层约3 cm上游和约5 cm下游的位置，密度归一化也主要使用中心测量。图5--7显示，流出受装置磁通核限制时，导引场在下游堆积并增加背压；观察到的速率下降大于所比较的局地预测。

作者明确表示，现有数据不能把局地Hall改变与全局几何／边界的贡献定量分开，仍需完整压力和离子流测量。因此这篇实验支持“导引场和边界共同影响重联”的判断，不支持对任意天体电流片直接使用同一个减速系数。

\subsection{2013长文：导引场压缩发生在比电流层更大的尺度}
\citet{L193}进一步展示约20 cm范围的导引场压缩，远大于约2 cm的片半宽。在反螺度操作中，局地导引场可达外加值约两倍；归一化率由无导引场附近约0.08，降到$B_g/B_{\rm rec}\simeq1.5$时约0.02。伴随流出区磁压增加，但实验并未把这一背压与局地Hall变化的贡献独立隔离。

这篇长文的Hall比较测点为上游约4 cm和下游约8 cm，与2012短文的测点不同，不能混作同一组剖面。测量位于X点外；由Hall项与局地电场相近不能推断Hall项直接提供X点的破冻结电场。图3的各面板采用不同色标，比较导引场前后的幅值应读数值刻度，不能凭颜色深浅。

配套Hall模拟采用$25\times12.5d_i$的Harris片，初始半厚$0.5d_i$、扰动0.05、$T_i=T_e$、总压绝热指数$5/3$、电阻0.005，流出周期而流入导体。其平板几何没有产生实验的全局导引场压缩；要取得较好局地比较，需使用\emph{实测局地增强后的}导引场。论文没有提供网格和步长，所以这不是一个已具完整输入的装置复现。

同螺度与反螺度操作的初始磁结构和全局演化也不同，不能把两类数据拼成只改变$B_g$的严格单参数扫描。真空环向场近$1/R$时，磁压力梯度与曲率张力本可相互抵消；讨论额外背压应针对被压缩的场及其几何，而非仅看到导引场就加一个未平衡的磁压。这是局地模型需要全局实验信息才能可靠解释的具体例子。

\section{如何把实验现象变成可证伪的流体检验}
对于Hall MHD，优先比较可测的四极场极性与形状、电子／离子流分离、片漂移和全局通量率。对于十矩双流体，进一步比较压力各向异性、非对角压力及欧姆定律分量。压力模型预测的新图样应与仪器能实际分辨的量联系；例如CGL的反极性四极图样是待实验检验的机制判据，不能因模拟中出现就写成已观测事实。\citep{Cassak2015}

相关实验／观测候选已单列在索引中。尚未完成全文核读的实验条目，不在本章据摘要评定是否“验证”某种流体模型。

\section{强密度不对称：MRX、MMS与PIC的跨尺度对照}
\citet{Yamada2018}比较MRX氦放电、2015年10月16日MMS磁层顶穿越和VPIC粒子模拟。\textbf{该文讨论双流体物理，但配套模拟不是双流体流体代码。}密度比约10、导引场很弱时，Hall四极场变成近双极结构；电子与离子的面内停滞点均偏离X点，且偏移量不同。MRX中电子停滞点偏离约 $5$--$10d_e$，离子约 $0.3$--$0.5d_i$。实验电子速度由磁场旋度、电荷密度和离子速度反演，而非直接逐粒子测速。

电子电磁能量转移集中在电子停滞点附近，垂直电功贡献显著大于平行贡献。离子受偏向低密度侧的势阱加速；MMS可看到离子反向流束，而较有碰撞的实验更接近热化分布。这些现象对标量压力模型提出直接挑战：即使平均流速和磁场相似，离子内能及分布函数也未必相同。

以输入磁能为1，原文表2给非对称MRX的出射磁能、离子获得能量、电子获得能量分别为0.44、0.31、0.25；配套PIC为0.43、0.32、0.25。这个积分比较针对选定近X点体积，不是宇宙中普遍的能量分配定律。MMS有限穿越轨迹不能独立测出相同的全局库存；为比较大尺度结构，该文还滤除了高于离子回旋频率的电场成分。因此不能用其图表排除高频波对能量交换的贡献。作者也未找到波幅与重联率之间的确定关系。

PIC采用质量比100、$40d_i\times20d_i$ 区域、$4200\times2100$ 网格，能量在 $2.5d_i\times3d_i$ 子域积分。方法段对周期方向和导体／反射方向的文字与单电流片坐标需进一步核对输入文件，报告保留这一复现疑点。实验、空间和模拟的形态一致是有价值的证据，但不足以直接判定某个五矩或十矩闭合已被验证。

\section{低Lundquist数也能有电阻性磁岛链}
Jara-Almonte等在MRX碰撞氩等离子体中发现电阻性亚离子电流片的磁岛链。尺度为
\begin{equation*}
 d_e\sim0.5\,\mathrm{mm}<\lambda_{\rm mfp}\sim2\,\mathrm{mm}
 <\delta\sim9\,\mathrm{mm}\ll d_i\sim20\,\mathrm{cm}.
\end{equation*}
片内电子温度约4 eV、密度约$2\times10^{20}\,\mathrm m^{-3}$，测得电阻接近Spitzer值；存在明显Hall四极磁场，但电子惯性可忽略。这是“碰撞性”与“双流体尺度”同时成立的例子。磁岛在$S<100$即出现，不能套用不可压缩单流体常见的$S\sim10^4$经验阈值。\citep{E448}

证据来自约400个磁探针：最佳径向分辨率6 mm、轴向3 cm。研究结合局地场插值的X/O点、Poincar\'e--Hopf绕数和磁场过境变号确认磁岛；仅凭积分磁通图识别小岛反而较不可靠。约500次磁岛形成事件中，超过四分之三伴随电场增强；200次放电的统计显示，平均一个磁岛对应约10\%的全球归一化重联率增量。没有磁岛时率已约0.1--0.2，因此磁岛不是这批实验“从慢到快”的唯一开关。

具体电子撕裂机制尚需含实际密度剖面、可压缩性、Hall场与有限长宽比的二维本征计算。实验还留下结构问题：为何$ d_i\gg\delta\gg d_e$的长薄层能存在，而未自动变为电子尺度开放X形？可压缩电阻Hall模型有机会解决其中一部分，但需复现实验驱动和尺度，不能只画相似磁岛。

\section{局域三维启动：感应场传播不等于总电场传播}
Katz等的VTF环形氩实验使用强导引场，密度约$2\times10^{18}\,\mathrm m^{-3}$、电子／离子温度约15／1 eV。在外部轴对称驱动下，重联仍先在局部环向位置启动。六个环向位置测得感应电场双向传播，正文报告速度约两倍Alfvén速度；加上静电部分后，总电场仍局域在启动附近。将摘要中的“近Alfvén”理解为精确等于$V_A$会丢失这一细节。\citep{Katz2010}

快重联与外区安全因子$q\simeq2$或3以及旋转电势模相伴；论文详析$q=2$。离子极化电流使局域下降的平行电流能够闭合。以实测势和磁场构造的半经验模型得到22微秒增长时间，与实验$20\pm6$微秒一致。但模型已使用拟合电势、实测感应场和二者共同增长的关系，不能视为从任意平衡无参数预测启动。浮动势近似、8 cm电流分辨率及边缘磁化的离子也需纳入比较。

\section{MAST：双峰密度的解释需要同时纳入Hall与装置几何}
Stanier等从总电流268 kA和估计磁绳半径构造MAST合并初值，并在模拟中按实验Thomson散射弦的位置$Z=15$ mm取径向密度剖面。环形内侧空间较小，磁绳靠近中心柱和重联喷流使内侧压缩；Hall电子平行流又形成四极样密度结构。两者叠加使诊断弦穿过两个高密度区，出现双峰；随合并结束，外侧峰减弱或消失，内侧峰保留。\citep{Stanier2013}

这是对实验\emph{形状和演化顺序}的解释，不是逐时定量重现：图中模拟每$20\tau_0\simeq5.8\,\mu$s取样，实验约每$0.1$ ms取样，峰值密度也不同；初始磁结构尚无法从实验完整反演。模型强制$T_i=T_e$，离子黏性和超电阻增强，故不能由这篇直接推断离子、电子的真实加热分配。它最有价值的验证策略，是先比较诊断真正看到的弦剖面，再追溯产生剖面的二维结构。

\section{2026年跨尺度对照：延伸双峰电流片与能量分配差异}
Yoo等比较MRX、2017年6月19日MMS磁尾穿越、2021年4月29日PSP日球电流片穿越及VPIC模拟。核心新结果是：分离线两侧的窄电流峰可以沿宏观排气区持续延伸，并伴随法向电势井、离子加速和密度增强。它把局地电子／离子解耦与宏观开放排气边界联系起来；\textbf{配套数值为PIC，不是Hall流体模拟}。\citep{L312}

MRX在离X线约$6d_i$处用高分辨探针拟合双Harris片，两个半厚约$0.29d_i$和$0.70d_i$。二维阵列本身分辨率不够，电流峰不能只靠其粗网格图确认。密度、离子速度剖面来自约250个出现清晰目标结构的重复放电；$T_i=T_e$则依据相近条件的氦放电作估计，而非这批氢放电逐点直接测量。

\begin{center}\small
\begin{tabular}{lrrl}\toprule
环境 & 势井估算（两侧） & 势井测量（两侧） & 主要限制\\\midrule
MRX &12 / 6 V&11 / 5 V&浮动势代理、密度非对称\\
MMS &6 / 5 kV&5 / 7 kV&法向速度估计、时间积分\\
PSP &120 / 60 V&未可靠测得&电场误差过大\\\bottomrule
\end{tabular}\end{center}
第一章式\eqref{eq:separatrix_potential}给出相应磁压差标度。MRX的数量接近，MMS则是量级相符，不能写成各环境都精确验证。MMS电流与$E\times B$漂移电流的吻合只在强法向电场的边界时窗成立；复杂多维区还有没有对应磁场台阶的电流峰。

能量分配也不同：扣除上游温度后，MMS离子焓流仍超过体动能通量，PSP则体动能通量较大。作者提出低杂漂移与相混合可能影响这种差异，但两个不同环境的个例不能单独隔离$\beta_e$的因果作用。实验中没有独立扫描出普适的$\beta$阈值。

大盒PIC为$200\times40d_i$、$4032\times800$网格、质量比100、$T_i/T_e=1$；小盒为$20\times20d_i$、$1280^2$、质量比1836和温比5。二者不是仅改变系统尺寸的对照。大盒内$\beta=0.012$时大岛破坏边界延伸，0.12／0.36时可局部延伸，1.2时保持到计算边界；这支持在\emph{受测参数}下的准稳态条件，没有给出所有导引场、质量比和三维条件的相图。

Hall／双流体下一步可直接检验电流双峰间距、厚度随下游距离的变化、势差与磁压关系，以及采用同样时间滤波的合成轨道。低杂波、束流与准确热化仍应由适当张量闭合和动理学对照验证。该文为Article in Press；本地保留了原文与透明评议文件，源数据Figshare接口在检索日返回404，故未声称完成原始数据复算。
